% GENERATED FILE. Edit canonical lessons, then run npm run generate:books.
% canonical-source-hash: fnv1a64:6f49761fe71774d0
% canonical-lessons: PA-C203-book-return, PA-C203-baare, PA-C203-baare-review, PA-C203-about-the-book, PA-C203-payoff

\chapter{About the Book}
\label{ch:pa-a2-baare}
\hlchaptermodality{203}

\begin{chapteropening}
\textbf{By the end of this chapter:} \emph{I can recognise and use the postposition for about in one bounded familiar-noun frame.}

Recognise and produce the about-postposition in one bounded familiar-noun frame through listening, speaking, reading, and writing.
\end{chapteropening}

\section[\texorpdfstring{\pa{ਕਿਤਾਬ} (kitaab)}{ਕਿਤਾਬ (kitaab)}]{Return to the familiar book noun}

\label{lesson:PA-C203-book-return}



\begin{quote}
\emph{Hear:} \emph{kitāb}; YOU SAY: \textbf{\pa{ਕਿਤਾਬ}}
\end{quote}

\begin{grammarlens}[title={Grammar Lens}]
Keep this familiar noun ready. A new relationship word will follow it only after the next lesson owns that word.
\end{grammarlens}

\subsection*{Your turn}
\begin{itemize}
\raggedright
  \item \emph{Hear:} \emph{book}; YOU SAY: \textbf{\pa{ਕਿਤਾਬ}}
  \item Say the familiar noun once, then leave a silent space after it.
\end{itemize}

\begin{tcolorbox}[breakable,colback=teal!4,colframe=teal!35!black,title={Before you move on}]
Produce the familiar noun from meaning. This opener owns no new word.
\end{tcolorbox}
\section[\texorpdfstring{\pa{ਬਾਰੇ} (baare)}{ਬਾਰੇ (baare)}]{One postposition for about}

\label{lesson:PA-C203-baare}



\begin{quote}
Retrieve the familiar noun for a book, then leave a silent space after it.
\end{quote}

\begin{grammarlens}[title={Grammar Lens: direct owner}]
\begin{quote}
\textbf{\pa{ਬਾਰੇ}} — \emph{bāre} — about; concerning
\end{quote}

Hear and say the whole word before reading it. Keep it as one whole postposition. Punjabi places a postposition after its noun; the next lesson will retrieve this new word before the familiar noun returns.
\end{grammarlens}

\begin{cousinweb}
Punjabi University Patiala explains the after-the-noun position. A dictionary entry records the “about / concerning” meaning and cites Bhai Maya Singh's 1895 dictionary. Punjab School Education Board Class 5 material independently uses the word in ordinary school prompts. Sources: \href{https://www.learnpunjabi.org/postposition.html}{Punjabi University Patiala}, \href{https://en.wiktionary.org/wiki/\%E0\%A8\%AC\%E0\%A8\%BE\%E0\%A8\%B0\%E0\%A9\%87}{dictionary entry}, and \href{https://static.pseb.ac.in/media/1654685607\_5th\%20Punjabi-second-language.pdf}{PSEB Class 5}.

The exact historical derivation of modern written \textbf{\pa{ਬਾਰੇ}} remains unknown in the cited sources. Its origin remains unknown here. This lesson therefore makes no etymological claim. Keep the present-day word whole.
\end{cousinweb}

\subsection*{Your turn}
\begin{itemize}
\raggedright
  \item \emph{Hear:} \emph{bāre}; YOU SAY: \textbf{\pa{ਬਾਰੇ}}
  \item \emph{Read it:} \textbf{\pa{ਬਾਰੇ}}
  \item \emph{Write it:} copy \textbf{\pa{ਬਾਰੇ}} once, checking the whole word
\end{itemize}

\begin{tcolorbox}[breakable,colback=teal!4,colframe=teal!35!black,title={Before you move on}]
Produce the postposition for “about / concerning” without looking.
\end{tcolorbox}
\section[\texorpdfstring{\pa{ਬਾਰੇ} (baare)}{ਬਾਰੇ (baare)}]{Retrieve the new postposition on its own}

\label{lesson:PA-C203-baare-review}



\begin{quote}
\emph{Hear:} \emph{bāre}; YOU SAY: \textbf{\pa{ਬਾਰੇ}}
\end{quote}

\subsection*{Your turn}
\begin{itemize}
\raggedright
  \item \emph{Hear:} \emph{about or concerning}; YOU SAY: \textbf{\pa{ਬਾਰੇ}}
  \item \emph{Read it:} \textbf{\pa{ਬਾਰੇ}}
  \item \emph{Write it:} cover and write \textbf{\pa{ਬਾਰੇ}}
\end{itemize}

\begin{grammarlens}[title={Grammar Lens}]
Keep one whole postposition and one bounded meaning here. A familiar noun will return only after this separate retrieval.
\end{grammarlens}

\begin{tcolorbox}[breakable,colback=teal!4,colframe=teal!35!black,title={Before you move on}]
Say and then write the word once from memory.
\end{tcolorbox}
\section[\texorpdfstring{\pa{ਕਿਤਾਬ} \pa{ਬਾਰੇ} (kitaab baare)}{ਕਿਤਾਬ ਬਾਰੇ (kitaab baare)}]{Put the postposition after one familiar noun}

\label{lesson:PA-C203-about-the-book}



\begin{quote}
Retrieve the familiar book noun, then the new postposition.
\end{quote}

\begin{grammarlens}[title={Grammar Lens}]
\begin{quote}
\textbf{\pa{ਕਿਤਾਬ} \pa{ਬਾਰੇ}} — \emph{kitāb bāre} — about the book
\end{quote}

Punjabi puts this postposition after the familiar noun: book first, then “about.” This chapter owns only that relationship word in this one gentle frame; the book noun remains old.
\end{grammarlens}

\subsection*{Your turn}
\begin{itemize}
\raggedright
  \item \emph{Hear:} \emph{kitāb bāre}; YOU SAY: \textbf{\pa{ਕਿਤਾਬ} \pa{ਬਾਰੇ}}
  \item \emph{Read it:} \textbf{\pa{ਕਿਤਾਬ} \pa{ਬਾਰੇ}}
  \item \emph{Write it:} cover and write \textbf{\pa{ਕਿਤਾਬ} \pa{ਬਾਰੇ}}
\end{itemize}

\begin{tcolorbox}[breakable,colback=teal!4,colframe=teal!35!black,title={Before you move on}]
Produce the complete “about the book” frame.
\end{tcolorbox}
\section[\texorpdfstring{\pa{ਕਿਤਾਬ} \pa{ਬਾਰੇ} (kitaab baare)}{ਕਿਤਾਬ ਬਾਰੇ (kitaab baare)}]{A bounded about-postposition payoff}

\label{lesson:PA-C203-payoff}



\begin{quote}
Retrieve the new postposition alone before rebuilding its familiar noun frame.
\end{quote}

\subsection*{Your turn: four skills}
\begin{itemize}
\raggedright
  \item \emph{Hear:} \emph{kitāb bāre}; YOU SAY: \textbf{\pa{ਕਿਤਾਬ} \pa{ਬਾਰੇ}}
  \item \emph{Read it:} \textbf{\pa{ਕਿਤਾਬ} \pa{ਬਾਰੇ}}
  \item \emph{Write it:} cover the line and write \textbf{\pa{ਕਿਤਾਬ} \pa{ਬਾਰੇ}}
\end{itemize}

\begin{tcolorbox}[breakable,colback=teal!4,colframe=teal!35!black,title={Before you move on}]
Say the new postposition, then produce the complete familiar-noun frame.

This is the twenty-sixth prerequisite tranche exposed by Punjabi A2 mock drafting. It does not unblock Punjabi A2 mock 1 by itself. It is not an A2 mock, not A2 readiness or attainment, and not external or human validation.
\end{tcolorbox}
